\section*{الفصل \ref{ch:b3:modules} --- المقاسات على حلقة مثاليات
رئيسية}

\begin{solution}{exo:b3:modules:1}
(a) إذا كان $\Q = \Z q_1 + \dots + \Z q_k$، فليكن $d$ مقامًا مشتركًا للأعداد $q_i$: عندئذٍ
تقع كل تركيبة في $\frac1d\Z$، لكن $\frac1{2d} \notin \frac1d\Z$. وهو تناقض.

(b) عديم \omterm{def:b3:modules:torsion}{الالتواء}: إذ يفرض $nq = 0$ مع $n \neq 0$ أن $q = 0$ في
$\Q$. وليس \omterm{def:b3:modules:free}{حرًّا}: إذ يحقق أي عددين ناطقين غير معدومين $\frac ab, \frac cd$
العلاقةَ غير التافهة $(bc)\frac ab - (ad)\frac cd =
0$، ومنه فإن أي أساس يضم عنصرًا واحدًا
على الأكثر؛ ولو كان $\Q \cong \Z$ لصار $\Q = \Z q$ دائريًا، لكن
$\frac q2 \notin \Z q$. (و $\Q \neq
0$.)

(c) تفترض \cref{thm:b3:modules:structure} \emph{التوليد المنتهي}، وهو ما ينفيه (a):
فلا تناقض --- بل إن $\Q$ يبيّن أن الفرضية لازمة في العبارة
«عديم \omterm{def:b3:modules:torsion}{الالتواء} $\Rightarrow$ حر».
\end{solution}

\begin{solution}{exo:b3:modules:2}
لدينا $360 = 2^3\cdot3^2\cdot5$. والتجزئات: للعدد $3$: $(3), (2,1),
(1,1,1)$؛ وللعدد $2$:
$(2), (1,1)$؛ وللعدد $1$: $(1)$. ومنه $3 \times 2
\times 1 = 6$ زمرة. ومن \omterm{thm:b3:modules:structure}{القواسم
الأوّلية} $\to$ \omterm{thm:b3:modules:smith}{العوامل الصامدة}:
\begin{center}
\begin{tabular}{l l}
$\Z/8 \times \Z/9 \times \Z/5$ & $\Z/360$\\
$\Z/8 \times \Z/3 \times \Z/3 \times \Z/5$ & $\Z/3 \times
\Z/120$\\
$\Z/4 \times \Z/2 \times \Z/9 \times \Z/5$ & $\Z/2 \times
\Z/180$\\
$\Z/4 \times \Z/2 \times \Z/3 \times \Z/3 \times \Z/5$ & $\Z/6
\times \Z/60$\\
$(\Z/2)^3 \times \Z/9 \times \Z/5$ & $\Z/2 \times \Z/2 \times
\Z/90$\\
$(\Z/2)^3 \times \Z/3 \times \Z/3 \times \Z/5$ & $\Z/2 \times
\Z/6 \times \Z/30$
\end{tabular}
\end{center}
(وللانتقال إلى \omterm{thm:b3:modules:smith}{العوامل الصامدة}: يجمع العامل الأكبر $d_s$ أعلى
قوة لكل عدد أولي، وهكذا نزولًا.) وأما ذات الرتبة $p^5$: فعددها
عدد تجزئات $5$، أي $7$.
\end{solution}

\begin{solution}{exo:b3:modules:3}
المصفوفة $B$: $D_1 = \gcd(2,4,6,8) = 2$؛ $D_2 = \abs{\det B} = \abs{16 -
24} = 8$. \omterm{thm:b3:modules:smith}{والعوامل الصامدة} $d_1 = 2$
و $d_2 = 8/2 = 4$: فصيغة سميث $\operatorname{diag}(2, 4)$، و $\Z^2/B\Z^2 \cong \Z/2\Z
\times \Z/4\Z$.

المصفوفة $C$: قطرية لكنها \emph{ليست} صيغة سميث (لأن
$2 \nmid 3$). لدينا $D_1 =
\gcd(2,3,12) = 1$؛ $D_2 = \gcd(2\cdot3,\, 2\cdot12,\, 3\cdot12) =
\gcd(6, 24, 36) = 6$؛ $D_3 = 72$. ومنه
$d = (1, 6, 12)$ و $\Z^3/C\Z^3 \cong \Z/6\Z\times\Z/12\Z$ --- وهذا متسق مع مبرهنة الباقي
الصيني: $\Z/2\times\Z/3\times\Z/12 \cong \Z/6\times\Z/12$.
\end{solution}

\begin{solution}{exo:b3:modules:4}
نكتب $B = Q\,\operatorname{diag}(d_1, \dots, d_n)\,P$ حيث $Q, P
\in GL_n(\Z)$ (\cref{thm:b3:modules:smith}؛ ولا يوجد $d_i$ معدوم لأن
$\det B \neq 0$). عندئذٍ $\Z^n/B\Z^n \cong \Z^n/
\operatorname{diag}(d)\Z^n = \prod_i \Z/d_i\Z$ (إذ يرسل التماثل المركّب $x \mapsto Q^{-1}x$
للزمرة $\Z^n$ المجموعةَ $B\Z^n$ على $\operatorname{diag}(d)P\Z^n = \operatorname{diag}(d)\Z^n$). وعدد عناصره $\prod \abs{d_i} = \abs{\det
\operatorname{diag}(d)} = \abs{\det B}$،
لأن $\det Q, \det P = \pm
1$. ومن أجل $L = \Z(2,0) + \Z(1,3)$: $B = \bigl(\begin{smallmatrix}2
& 1\\ 0 & 3\end{smallmatrix}\bigr)$، $D_1 = 1$، $D_2 = 6$: ومنه
$\Z^2/L
\cong \Z/6\Z$، وعدد عناصره $\abs{\det B} = 6$.
\end{solution}

\begin{solution}{exo:b3:modules:5}
(a) \omterm{def:b3:modules:module}{مقاس} جزئي: وقد تم ذلك في \cref{def:b3:modules:torsion}. وإذا كان $a(x +
T(M)) = 0$ في
$M/T(M)$ مع $a \neq 0$، فإن $ax \in T(M)$: أي $bax
= 0$ من أجل
$b \neq 0$ ما، ويكون $ba \neq 0$ (لأن الحلقة تامة)، ومنه $x \in
T(M)$:
أي إن الصنف معدوم. إذن $M/T(M)$ عديم \omterm{def:b3:modules:torsion}{الالتواء}.

(b) \omterm{def:b3:modules:module}{المقاس} $(X, Y) \subseteq K[X,Y]$ عديم \omterm{def:b3:modules:torsion}{الالتواء} (لأنه \omterm{def:b3:modules:module}{مقاس} جزئي من الحلقة
التامة $K[X,Y]$ الفاعلة على نفسها). ولو كان \omterm{def:b3:modules:free}{حرًّا}: فإن أي
عنصرين $P, Q$ يحققان $Q\cdot P - P \cdot Q = 0$، وهي علاقة غير تافهة حين يكون
$P \ne Q$ غير معدومين، ومنه فإن أي أساس يضم عنصرًا واحدًا:
أي $(X, Y) = (P)$ رئيسي --- وهذا يناقض \cref{exo:b3:rings:6}(a). فهو إذن
عديم \omterm{def:b3:modules:torsion}{الالتواء} ومولَّد بعدد منتهٍ (لأن $X, Y$ يولّدانه)، ومع ذلك
ليس \omterm{def:b3:modules:free}{حرًّا}: فعلى الحلقة $K[X,Y]$ التي ليست \omterm{def:b3:rings:pidufd}{حلقة مثاليات رئيسية}،
تخفق مبرهنة البنية.
\end{solution}

\begin{solution}{exo:b3:modules:6}
(a) إذا كان $\Z = 2\Z \oplus C$، فإن الإسقاط $\Z \to \Z/2\Z$ يُقصَر إلى
تماثل $C \cong \Z/2\Z$: أي إن $C$ ستكون زمرة جزئية من $\Z$
يحقق عنصرها غير المعدوم $x$ العلاقةَ $2x \in C
\cap 2\Z = 0$. لكن $\Z$ عديم
\omterm{def:b3:modules:torsion}{الالتواء}: إذن $C = 0$، وهذا يفرض $\Z =
2\Z$ --- وهو خطأ.

(b) \omterm{def:b3:modules:module}{المقاس} $A^n/M$ مولَّد بعدد منتهٍ وعديم \omterm{def:b3:modules:torsion}{الالتواء}، ومنه فهو حر
(\cref{thm:b3:modules:structure}): أي $A^n/M \cong A^r$ بأساس $f_1, \dots, f_r$. نختار سوابق
$y_i \in A^n$ للعناصر $f_i$ ونضع $F = Ay_1 + \dots + Ay_r$. ولكل $x \in A^n$ لدينا
$\pi(x)
= \sum a_if_i$، ومنه $x - \sum a_iy_i \in M$: أي $A^n = M + F$. وإذا كان $\sum a_iy_i \in M$، فإن
تطبيق $\pi$ يعطي $\sum a_if_i = 0$، ومنه فإن جميع المعاملات $a_i = 0$ (لأن
العائلة أساس): أي $M \cap F = 0$. إذن $A^n = M \oplus
F$.
\end{solution}

\begin{solution}{exo:b3:modules:7}
كان اختزال \cref{exo:b3:modules:3} فعّالًا: إذ
\[
L = \begin{pmatrix} 1 & 0\\ -3 & 1\end{pmatrix},
\qquad
R = \begin{pmatrix} 1 & 2\\ 0 & -1 \end{pmatrix},
\qquad
LBR = \begin{pmatrix} 2 & 0\\ 0 & 4\end{pmatrix}
\]
(بالعملية على الأسطر $R_2 \leftarrow R_2 - 3R_1$، والعمليتين على الأعمدة $C_2 \leftarrow C_2 - 2C_1$ ثم
$C_2 \leftarrow -C_2$). وبوضع $y = R^{-1}x$ (وهو تقابل للمجموعة
$(\Z/20\Z)^2$، لأن $R$ قابلة للقلب على $\Z$)، تكافئ الجملة
$Bx \equiv b \pmod{20}$ الجملةَ
\[
2y_1 \equiv b_1, \qquad 4y_2 \equiv -3b_1 + b_2 \pmod{20}.
\]
وتكون المطابقة $ky \equiv c \pmod{20}$ \omterm{def:b3:groups:derived}{قابلة للحل} إذا وفقط إذا كان $\gcd(k, 20) \mid c$:
فالحلول موجودة إذا وفقط إذا كان $2 \mid b_1$ و $4
\mid b_2 - 3b_1$، أي
\emph{$b_1$ زوجي و $b_2 \equiv 3b_1
\pmod 4$}. وحين تكون \omterm{def:b3:groups:derived}{قابلة للحل} يكون عدد
الحلول $2 \times 4 = 8$ بترديد $20$.
\end{solution}

\begin{solution}{exo:b3:modules:8}
(a) التشاكل $u$ المعدوم القوة يحقق $\mu_u = X^k$؛ وتكون قواسمه
الأوّلية $X^{k_1} \geq \dots$، أي كتلة جوردان $J_{k_i}(0)$ واحدة لكل جزء من
تجزئة العدد $4$:
\begin{center}
\begin{tabular}{l l l}
التجزئة & صيغة جوردان & \omterm{thm:b3:modules:smith}{العوامل الصامدة}\\
$(4)$ & $J_4$ & $X^4$\\
$(3,1)$ & $J_3 \oplus J_1$ & $X,\ X^3$\\
$(2,2)$ & $J_2 \oplus J_2$ & $X^2,\ X^2$\\
$(2,1,1)$ & $J_2 \oplus J_1 \oplus J_1$ & $X,\ X,\ X^2$\\
$(1,1,1,1)$ & $0$ & $X, X, X, X$
\end{tabular}
\end{center}
(b) نأخذ $u = J_2\oplus J_2$ و $v = J_2 \oplus J_1 \oplus J_1$: فلكلتيهما $\chi = X^4$ و $\mu = X^2$،
لكن عواملهما الصامدة مختلفة --- ومنه فهما غير متشابهتين
(\cref{thm:b3:modules:frobenius})؛ ويمكن أيضًا مقارنة الرتبتين: $\operatorname{rk} u = 2 \neq 1 =
\operatorname{rk} v$. أما من أجل
$n \leq 3$: فإن $\chi$ و $\mu$ يحدّدان، من أجل كل قيمة ذاتية
$\lambda$ (على حقل انشطار)، القياسَ الكلي $m_\lambda \leq 3$ لكتل
$\lambda$ وأكبرَ كتلة $r_\lambda$؛ وتتحدّد تجزئة العدد
$m \leq 3$ بأكبر أجزائها (فحالة $m = 3$ و $r = 2$ تفرض $(2,1)$،
وهكذا). ومنه تتطابق \omterm{thm:b3:modules:structure}{القواسم الأوّلية}، وتُنزِل \cref{cor:b3:modules:descent} التشابهَ
إلى الحقل الأساسي.
\end{solution}

\begin{solution}{exo:b3:modules:9}
(أ)$\Rightarrow$(ب): إذا كان $V = K[u]x$، فإن $V \cong
K[X]/\operatorname{Ann}(x)$ و $\operatorname{Ann}(x) = (\mu_u)$
(لأن كثير حدود يُبيد $x$ إذا وفقط إذا أباد $V = K[u]x$ كلها، بما
أن $P(u)Q(u)x = Q(u)P(u)x$). ومنه $n = \dim V = \deg \mu_u$؛ وبما أن $\mu_u \mid \chi_u$ و $\deg\chi_u = n$،
تعطي الواحدية $\mu_u =
\chi_u$.

(ب)$\Rightarrow$(ج): لدينا $\deg\chi_u = \sum_i \deg P_i$ و $\mu_u = P_s$
(\cref{thm:b3:modules:frobenius})؛ ويفرض تساوي الدرجتين أن $s = 1$.

(ج)$\Rightarrow$(أ): إن $V \cong K[X]/(P_1)$ دائري، ويولّده سابق $\bar 1$.

\omterm{def:b3:modules:companion}{والمصفوفة المرافقة} هي حالة $V = K[X]/(P)$ نفسها: فالعنصر
$x = \bar
1$، أي $e_1$، دائري. ومن أجل مصفوفة قطرية $\operatorname{diag}(\lambda_1, \dots, \lambda_n)$: لدينا
$\chi = \prod
(X - \lambda_i)$ و $\mu = \prod_{\lambda \text{ متمايزة}} (X -
\lambda)$؛ وهما متساويان إذا وفقط إذا كانت العناصر
$\lambda_i$ متمايزة مثنى مثنى: أي إن لمصفوفة قطرية متجهةً
دائرية إذا وفقط إذا كانت عناصر قطرها متمايزة مثنى مثنى (وعندئذٍ
تصلح $x = (1, \dots, 1)$: بمصفوفة فاندرموند).
\end{solution}

\begin{solution}{exo:b3:modules:10}
سميث: $M = Q\operatorname{diag}(d_1,\dots,d_n)P$؛ وتعطي \cref{exo:b3:modules:4} أن $[\Z^n : M\Z^n] = \prod\abs{d_i} =
\abs{\det M}$. وإذا كان
$\det M = \pm 1$: فإن صيغة \omterm{def:b3:modules:companion}{المصفوفة المرافقة} $M^{-1} = (\det M)^{-1}\,{}^{t}\!\operatorname{com}(M)$ ذات عناصر
صحيحة، ومنه $M \in GL_n(\Z)$؛ وبالعكس، يعطي $MM^{-1} = I$ أن
$\det M \cdot \det M^{-1} = 1$ في $\Z$، ومنه $\det M = \pm1$.

الزمر الجزئية ذات الدليل $m$ في $\Z^2$: لمثل هذه الزمرة $L$
رتبةٌ تساوي $2$ (لأن دليلها منتهٍ) وأساسٌ وحيد في \emph{صيغة
هيرميت الناظمية} $\bigl(\begin{smallmatrix} a & b\\ 0 &
d\end{smallmatrix}\bigr)$: إذ يتميّز $d$ بالخاصية $L \cap (\{0\}
\times \Z) = \{0\} \times d\Z$،
ويتميّز $a$ بالعلاقة $\pi_1(L) = a\Z$ (الإحداثيات الأولى)، ثم
يكون $b$ وحيدًا بترديد $d$؛ وننظّم $a, d
> 0$ و $0 \leq b < d$.
والدليل هو $ad = m$. وأما الإحصاء: فمن أجل كل قاسم $d \mid m$
(حيث $a = m/d$) توجد $d$ اختيارات للعنصر $b$: فالمجموع $\sum_{d \mid m} d = \sigma_1(m)$.
\end{solution}

\begin{solution}{exo:b3:modules:11}
(a) بمبرهنة البنية، $G \cong \prod_i\Z/d_i\Z$، وتنفصل المعادلة $x^k = e$ إحداثيًا.
وفي $\Z/d\Z$: للمعادلة $kx \equiv 0
\pmod d$ عدد حلول قدره $\gcd(k, d)$ بالضبط (إذ
يجب أن يكون $x$ مضاعفًا للعدد $d/\gcd(k,d)$، وعدد هذه المضاعفات
$\gcd(k, d)$). ثم نضرب على العوامل.

(b) إذا كانت $G = \Z/d_s\Z$ دائرية ($s = 1$)، فالعدد $\gcd(k, d_s) \leq k$.
وإذا كان $s \geq 2$: نأخذ $k = d_1$؛ فيكون العدد $\prod_i\gcd(d_1, d_i) = d_1^{\,s} > d_1$ (لأن
كل $\gcd$ يساوي $d_1$ بسلسلة القابلية للقسمة): أي إن
للمعادلة $x^{d_1} = e$ أكثر من $d_1$ حل.

(c) في حقل، لكثير الحدود $X^k - 1$ عدد جذور لا يتجاوز $k$
(\cref{ch:b3:rings}: فكثير حدود غير معدوم من الدرجة $k$ على حلقة تامة)،
ومنه فإن كل زمرة جزئية منتهية $G \leq K^\times$ تحقق محك (b): أي إن $G$
دائرية --- وهو البرهان البنيوي من سطر واحد على دائرية $\mathbb F_q^\times$،
مكمّلًا برهان الإحصاء في \cref{ch:b3:galois}.
\end{solution}

\begin{solution}{exo:b3:modules:12}
(a) إذا كان $MN = I$ حيث $N$ صحيحة: فإن $\det M\det N = 1$ وكلاهما صحيح،
ومنه $\det M = \pm1$. وبالعكس، إذا كان $\det M =
\pm1$، فإن صيغة
المرافقات $M^{-1} =
\frac1{\det M}\,{}^t\!\operatorname{com}(M)$ ذات عناصر صحيحة.

(b) الضرب من اليسار في $E^{-k}$ يطرح $k$ مرة السطرَ $2$ من
السطر $1$؛ والضرب في $F^{-k}$ يطرح $k$ مرة السطرَ $1$ من السطر
$2$. وإذا أُعطيت $M = \bigl(\begin{smallmatrix}a & *\\ c &
*\end{smallmatrix}\bigr) \in SL_2(\Z)$، فإن العمود الأول $(a, c)$ متجهة
أحادية النمط ($\gcd(a, c) = 1$: لأنه يقسم $\det M = 1$). نشغّل
خوارزمية إقليدس على $(a, c)$ بهذه العمليات على الأسطر: فبعد عدد
منتهٍ من الخطوات يصير العمود $(\pm1, 0)$. وتصير المصفوفة عندئذٍ
$\pm
\bigl(\begin{smallmatrix}1 & b\\ 0 & 1\end{smallmatrix}
\bigr) = \pm E^{b}$ (لأن المحدد بقي $1$). ويبقى أن نكتب $-I$ بدلالة
المولِّدات: $(EF^{-1}E)^2 =
\bigl(\begin{smallmatrix}0 & 1\\ -1 &
0\end{smallmatrix}\bigr)^2 = -I$ (تحقق من مربّع مصفوفة الدوران). وبفكّ
الخطوات، تكون $M$ كلمةً في $E^{\pm1},
F^{\pm1}$.

(c) تستعمل خوارزمية سميث (\cref{met:b3:modules:smith}) عمليات على الأسطر
والأعمدة من هذا النوع بالضبط (مع التبديلات وتغييرات الإشارة،
وهي نفسها جداءات لعمليات أوّلية إلى حدود إشارة المحدد): فعلى
$\Z$، تُكتب كل مصفوفة على الصورة $U\,D\,V$ حيث $U, V$ جداءات
مصفوفات أوّلية و $D$ صيغة سميث --- والسؤال (b) هو الحالة
$2\times2$ ذات المحدد $1$ من الحقيقة العامة القائلة إن مصفوفات
النمط $E$ تولّد $SL_n(\Z)$.
\end{solution}

\begin{solution}{pb:b3:modules:1}
\textbf{1.} التطبيق $\varphi$ جمعي وخطي على $K[X]$: $\varphi(X
\cdot (Q_i)_i) = \sum_i (XQ_i)(M)e_i = M\sum_i Q_i(M)e_i = X
\cdot \varphi\bigl((Q_i)_i\bigr)$،
لأن بنية \omterm{def:b3:modules:module}{المقاس} $V_M$ هي $X \cdot v = Mv$. وهو شامل: لأن
المتجهات الثابتة تعطي $K^n$ كلها. والعمود $j$ من $XI - M$ هو
$Xe_j - \sum_i
m_{ij}e_i$، وصورته $Me_j - \sum_i m_{ij}e_i = 0$.

\textbf{2.} بترديد الأعمدة، $Xe_j \equiv \sum_i m_{ij}e_i$: فكل متجهة من كثيرات الحدود
تختزل، بالتراجع على أعلى درجة، إلى متجهة \emph{ثابتة}
$c \in K^n$. وإذا كانت المتجهة الأصلية في $\ker\varphi$، فإن
$\varphi(c) = c = 0$ (لأن $\varphi$ على الثوابت هو التطابق $K^n = V_M$)،
ومنه تقع المتجهة في الفضاء المولَّد بالأعمدة: أي $\ker\varphi = (XI_n -
M)K[X]^n$. ومنه
$V_M \cong K[X]^n/(XI - M)K[X]^n$، وتعطي صيغة سميث على $K[X]$ (وجميع عواملها الصامدة غير
معدومة لأن جداءها $\det(XI - M) = \chi_M$) أن $V_M \cong \bigoplus_i
K[X]/(f_i)$: فكثيرات الحدود $f_i$
غير الثابتة هي \omterm{thm:b3:modules:frobenius}{صوامد التشابه}، ويمكن حسابها بوصفها خوارج قواسم
مشتركة كبرى للمحدّدات الجزئية (\cref{thm:b3:modules:smith}).

\textbf{3.} من أجل $\lambda I_n$: تكون $XI - \lambda I$ في صيغة
سميث أصلًا: فالصوامد $(X - \lambda, \dots, X - \lambda)$، وعددها $n$. ومن أجل مصفوفة قطرية
عناصر قطرها متمايزة: $V \cong \bigoplus_i K[X]/(X -
\lambda_i)$ بموديلات متكاملة مثنى مثنى، ومنه
تضغط مبرهنة الباقي الصيني ذلك إلى \omterm{def:b3:modules:module}{المقاس} الدائري الوحيد $K[X]/\bigl(\prod_i(X - \lambda_i)\bigr)$:
أي صامد واحد هو $\chi$. ومن أجل $J_n(0)$: يفرض $\mu = X^n = \chi$ صامدًا
وحيدًا هو $X^n$. ومن أجل $\operatorname{diag}(J_2(0), J_1(0))$: \omterm{thm:b3:modules:structure}{القواسم الأوّلية} $X^2, X$:
فالصوامد $P_1 = X \mid P_2 =
X^2$.

\textbf{4.} يرسل $P \mapsto P(u)$ الحلقةَ $K[X]$ على $K[u]$،
ونواته $(\mu_u)$ بحكم تعريف كثير الحدود الأدنى: ومنه $K[u]
\cong K[X]/(\mu_u)$،
وبُعده $\deg\mu_u = \deg P_s = n_s$.

\textbf{5.} تفرض الخطية على $K[X]$ أن $f(\bar Q) = f(Q\cdot\bar 1)
= Q\,f(\bar 1)$. ويحقق الصنف
$\bar 1$ العلاقة $P \bar 1 = 0$، ومنه فإن $P f(\bar 1) = 0$ شرط لازم.
وبالعكس، إذا كان $P\bar c = 0$، فإن $f(\bar Q) = Q\bar c$ معرَّف تعريفًا
سليمًا (لأن $Q \equiv Q' \bmod P
\Rightarrow (Q - Q')\bar c \in$ مضاعفات $P\bar c = 0$) وخطي على $K[X]$.

\textbf{6.} ليكن $g = \gcd(P, Q)$ و $P = gP'$ و $Q = gQ'$ حيث
$\gcd(P', Q') = 1$. في $K[X]/(Q)$: $P\bar c = 0 \iff Q \mid Pc
\iff Q' \mid P'c \iff Q' \mid c$ (إقليدس، $\gcd(P', Q') = 1$). ومنه فإن
العناصر المقبولة $\bar c$ تشكّل \omterm{def:b3:modules:module}{المقاس} الجزئي المولَّد
بالعنصر $\bar{Q'}$، ومُبيده $\{R : Q \mid RQ'\} = (g)$: أي إن ذلك \omterm{def:b3:modules:module}{المقاس} الجزئي
$\cong K[X]/(g)$. ومع السؤال 5،
$\operatorname{Hom}(K[X]/(P), K[X]/(Q)) \cong K[X]/(\gcd(P,Q))$
وبُعده $\deg\gcd(P,Q)$.

\textbf{7.} يتبادل $v$ مع $u$ إذا وفقط إذا تبادل مع كل $P(u)$،
أي إذا وفقط إذا كان $v$ خطيًا على $K[X]$: أي $\mathcal C(u) =
\operatorname{End}_{K[X]}(V)$. وبكتابة
تشاكلات $V =
\bigoplus_j V_j$ على صورة مصفوفات $(f_{ij})$، أي $f_{ij} \in
\operatorname{Hom}(V_j, V_i)$ (بالتركيب
مع الانغمارات والإسقاطات)، يعطي السؤال 6
\[
\dim \mathcal C(u) = \sum_{i,j} \deg\gcd(P_i, P_j)
= \sum_{i,j} n_{\min(i,j)}
= \sum_{k=1}^{s} \bigl(2(s - k) + 1\bigr)\,n_k ,
\]
باستعمال سلسلة القابلية للقسمة ($\gcd(P_i, P_j) =
P_{\min(i,j)}$) وباستعمال، في الخطوة
الأخيرة، أن $\min(i,j) = k$ يتحقق من أجل $2(s - k) + 1$ زوجًا
$(i, j)$ بالضبط.

\textbf{8.} بما أن $2(s-k)+1 \geq 1$، فإن $\dim\mathcal C(u) \geq
\sum_k n_k = n$، مع التساوي إذا وفقط إذا
كان $s = 1$، أي إذا وفقط إذا كان $u$ دائريًا (\cref{exo:b3:modules:9}). ومن
أجل $u = \lambda\,\mathrm{id}$: يكون $s = n$ وجميع $n_k = 1$: $\dim = \sum_{k=1}^n (2(n-k)+1) = n^2$ --- وهذا
صحيح، لأن $\mathcal C(\lambda\,\mathrm{id}) = \mathcal
L(V)$.

\textbf{9.} بالصيغة: الصوامد $(X, X^2)$، ومنه $s = 2$ و $n_1 =
1$
و $n_2 = 2$: أي $\dim = 3\cdot1 + 1\cdot2 = 5$. ومباشرةً: في الأساس $(e_1, e_2, e_3)$ حيث
$u e_2 = e_1$ و $ue_1 = ue_3 = 0$، يعطي كتابة $Au = uA$ من أجل
$A = (a_{ij})$ الشروطَ $a_{21} = a_{23} = a_{31} = 0$ و $a_{11} = a_{22}$: أي خمسة وسائط حرة
$a_{11}{=}a_{22}, a_{12}, a_{13}, a_{32}, a_{33}$.

\textbf{10.} كل كثير حدود $P(u)$ يتبادل مع كل ما يتبادل مع $u$
(لأنه مجموع قوى للتشاكل $u$): ومنه $K[u]
\subseteq \mathcal C(\mathcal C(u))$. و $u \in \mathcal C(u)$، ومنه فإن أي
$w \in \mathcal C(\mathcal C(u))$ يتبادل مع $u$.

\textbf{11.} ليكن $V = K[u]x$ وليكن $v \in \mathcal C(u)$. نكتب $v(x) = P(u)x$
(بالدائرية). ومن أجل $y = Q(u)x$ كيفي: $v(y) =
vQ(u)x = Q(u)v(x) = Q(u)P(u)x = P(u)y$. إذن $v = P(u)$:
أي $\mathcal
C(u) = K[u]$. ومنه $\mathcal C(\mathcal C(u)) = \mathcal
C(K[u]) = \mathcal C(u) = K[u]$ (لأن التبادل مع $K[u]$ كلها هو نفسه
التبادل مع $u$). فالمبرهنة صحيحة في الحالة الدائرية.

\textbf{12.} التطبيق $\pi_i$ خطي على $K[X]$ (لأن التفكيك مجموع
مباشر لمقاسات جزئية)، ومنه $\pi_i \in \mathcal C(u)$، ويتبادل $w$ معه: $w(V_i) = w\pi_i(V) = \pi_i w(V) \subseteq
V_i$.
والقصر $w_i = w\restriction_{V_i}$ يتبادل مع $u_i = u\restriction_{V_i}$ الدائري (السؤال 10)، و $w_i
\in \mathcal C(u_i) = K[u_i]$
(السؤال 11): أي $w_i = Q_i(u_i) =
Q_i(u)\restriction_{V_i}$ من أجل $Q_i \in K[X]$ ما.

\textbf{13.} سلامة تعريف $\eta_{ij}(R(u)x_j) = R(u)x_i$: إذا كان $R(u)x_j = 0$ فإن
$P_j \mid R$، ويعطي $P_i \mid P_j$ أن $P_i \mid R$، ومنه
$R(u)x_i = 0$ ($\operatorname{Ann}(x_i) =
(P_i)$). وعندئذٍ يكون $\eta_{ij}$ خطيًا على
$K[X]$ بحكم البناء، ويكون $\tilde\eta_{ij} = \eta_{ij}\pi_j$ تركيبًا لتشاكلات $V \to V$ على
$K[X]$: أي $\tilde\eta_{ij} \in \mathcal
C(u)$.

\textbf{14.} نقيّم $w\tilde\eta_{ij} = \tilde\eta_{ij}w$ عند $x_j$: فالطرف الأيسر هو $w(x_i) = Q_i(u)x_i$؛
والطرف الأيمن هو $\eta_{ij}\bigl(Q_j(u)x_j\bigr) = Q_j(u)x_i$. ومنه $(Q_i -
Q_j)(u)\,x_i = 0$: أي $P_i \mid Q_i - Q_j$ من أجل كل
$i \leq j$. وعلى وجه الخصوص، مع $Q = Q_s$: $Q \equiv Q_i \pmod{P_i}$، ومنه
يتوافق $Q(u)$ و $Q_i(u)$ على $V_i$ (الذي يبيده $P_i(u)$).
ومن ثَمّ $w
= Q(u)$ على كل $V_i$، ومنه على $V$: أي $\mathcal C(\mathcal C(u))
\subseteq K[u]$، ومع
السؤال 10، $\mathcal C(\mathcal
C(u)) = K[u]$.

\textbf{15.} العبارة الأولى هي الأسئلة من 10 إلى 14 (أو، من أجل
$u$ دائري، السؤال 11 وحده). ومن أجل $u = \mathrm{id}$، $n \geq
2$: يكون
بُعد $\mathcal C(u) = \mathcal L(V)$ مساويًا $n^2$، بينما $\dim K[u] = \deg\mu_u = 1$. ومنه فإن $\mathcal C(u) = K[u]$
يخفق إخفاقًا شديدًا؛ ومع ذلك تصح مبرهنة المُبادِل المضاعف
($\mathcal C(\mathcal L(V)) = K\,\mathrm{id} = K[u]$: فمركز جبر المصفوفات هو المصفوفات السلَّمية). فالدائرية
هي ما يجعل المُبادِل \emph{المفرد} كثيرَ حدود أصلًا؛ أما
المُبادِل \emph{المضاعف} فهو كثير حدود دائمًا.

\textbf{16.} إذا كان $P(XI - M)Q = S$ اختزالًا لسميث ($P, Q$
قابلتان للقلب على $K[X]$)، فإن النقل يعطي ${}^tQ\,(XI -
{}^tM)\,{}^tP = {}^tS = S$: أي صيغة سميث
نفسها، ومنه فإن للمصفوفتين $XI - M$ و $XI - {}^tM$ العواملَ
الصامدة نفسها، أي إن للمصفوفتين $M$ و ${}^tM$ صوامدَ التشابه
نفسها (\cref{cor:b3:modules:descent}): إذن هما متشابهتان.

\textbf{17.} \omterm{thm:b3:modules:frobenius}{صوامد التشابه} للمصفوفة $M$ على $L$ هي \omterm{thm:b3:modules:smith}{العوامل
الصامدة} للمصفوفة $XI - M$ في $L[X]$. وأي اختزال لسميث للمصفوفة
$XI - M$ على $K[X]$ --- بمصفوفتين $P, Q$ قابلتين للقلب على
$K[X]$، وقطري مع سلسلة القابلية للقسمة --- هو \emph{أيضًا}
اختزال سليم لسميث على $L[X]$ (لأن $P, Q$ تبقيان قابلتين للقلب:
فمحدداهما ثابتان غير معدومين)، \omterm{thm:b3:modules:smith}{والعوامل الصامدة} الواحدية وحيدة:
ومنه تتطابق \omterm{thm:b3:modules:smith}{العوامل الصامدة} المحسوبة على $K$ وعلى $L$. إذن
$M \sim_L N$ إذا وفقط إذا كان لهما \omterm{thm:b3:modules:smith}{العوامل الصامدة} نفسها إذا
وفقط إذا كان $M \sim_K N$. وعلى وجه الخصوص، تكون المصفوفتان
الحقيقيتان المقترنتان على $\C$ مقترنتين على $\R$ --- وهي عبارة
كثيرًا ما تُبرهَن تحليليًا (بتخصيص مصفوفة قابلة للقلب $P + \iu
Q$)،
وهي هنا بنيوية.

\textbf{18.} نفكّك $u = \bigoplus J_{m_t}(0)$ إلى كتل جوردان معدومة القوة. وفي كتلة
واحدة قياسها $m$، لدينا $\operatorname{rk}J_m^k = \max(m - k, 0)$، ومنه $\operatorname{rk}J_m^{k-1} - \operatorname{rk}J_m^k = 1$ إذا كان $m
\geq k$،
و $0$ فيما عدا ذلك. وبالجمع على الكتل: $\operatorname{rk}u^{k-1} - \operatorname{rk}u^k =
\#\{t : m_t \geq k\}$. ومن ثَمّ تحدّد
متتالية الرتب المجموعةَ المتعدّدة $(m_t)$ --- وهي تجزئة للعدد
$n$ --- وبالعكس تتحقق كل تجزئة: فأصناف معدومات القوة $\leftrightarrow$
تجزئات $n$، على كل حقل. ومن أجل $M_5$: عدد الأصناف $p(5) = 7$
($5$؛ $4{+}1$؛ $3{+}2$؛ $3{+}1{+}1$؛ $2{+}2{+}1$؛
$2{+}1{+}1{+}1$؛ $1^5$).

\textbf{19.} (a) على $\Q$ (أو على أي حقل مميّزه $0$)، لدينا
$D^n = 0$ و $D^{n-1}(X^{n-1}) = (n-1)!\, \neq 0$: أي إن $D$ معدوم القوة من الدليل $n$ على
فضاء بُعده $n$، ومنه فهو دائري بصامد وحيد $X^n$ (إذ يولّد
$x = X^{n-1}$: لأن مشتقاته المتكررة تولّد الفضاء). (b) على
$\mathbb F_p$ مع $p
< n$: لدينا $D^p = 0$، لأن المشتقة النونية
من الرتبة $p$ لكل حد أحادي $X^m$ تحمل العامل $m(m-1)\cdots(m-p+1)$، وهو جداء
$p$ عددًا صحيحًا متتاليًا، ومنه $\equiv 0 \bmod p$. نكتب $n = ap + r$ حيث
$0 \leq r < p$. عندئذٍ يتولّد $\ker D^k$ بالحدود الأحادية $X^m$
التي تحقق $D^kX^m = 0$؛ وبإحصاء الأُسس $m < n$ حسب بواقيها
بترديد $p$: $\dim\ker D^k = ak + \min(r,
k)$ من أجل $0 \leq k \leq p$، ومنه $\operatorname{rk}D^{k-1} -
\operatorname{rk}D^k = \dim\ker D^k - \dim\ker D^{k-1} = a +
\mathbf 1_{k \leq r}$. وحسب السؤال
18، تضم التجزئة $a$ كتلة قياسها $p$ بالضبط و (إذا كان $r > 0$)
كتلةً واحدة قياسها $r$: أي إن \omterm{thm:b3:modules:frobenius}{صوامد التشابه} $X^r \mid X^p \mid \dots \mid X^p$. فالمميّز
يغيّر الصيغة الناظمية لأشهر مؤثر في الرياضيات.

\textbf{20.} للمصفوفة $M \in GL_2$ عدد صوامد $s \in \{1, 2\}$.
فإذا كان $s = 2$: $P_1 = P_2 = X - a$ (مع $a \neq 0$ للقابلية للقلب)، أي
$M = aI$ مركزية. وإذا كان $s = 1$: تكون $M$ دائرية وكثيرُ حدودها
المميّز $=$ الأدنى $\chi$ من الدرجة $2$، وتوافق الأصنافُ
كثيراتِ الحدود $\chi$ الممكنة مع $\chi(0) \neq 0$: أي $\chi$
منشطر بجذرين متمايزين $a \neq
b$ (\omterm{def:b3:modules:companion}{والمصفوفة المرافقة} $\sim$ قطرية)؛
أو $\chi = (X - a)^2$ (المرافقة، غير نصف بسيطة)؛ أو $\chi$ غير قابل
للاختزال. ونمط واحد بالضبط لكلٍّ --- لأن \omterm{thm:b3:modules:smith}{العوامل الصامدة} صامدٌ
تام.

\textbf{21.} المركزية: $q - 1$ اختيارًا للعنصر $a$. والقيم
الذاتية المنشطرة المتمايزة: أزواج غير مرتّبة $\{a, b\} \subseteq
\mathbb F_q^\times$ مع
$a \neq b$: أي $\binom{q-1}2$ صنفًا. والأدنى $(X-a)^2$: أي
$q - 1$ صنفًا. وكثيرات الحدود من الدرجة الثانية غير القابلة
للاختزال ذات الحد الثابت غير المعدوم: جميعها صالحة (لأن جذورها
غير معدومة)، وعددها $\frac{q^2 - q}2$ من كثيرات الحدود الواحدية غير
القابلة للاختزال من الدرجة الثانية (أي $q^2$ من كثيرات الحدود
الواحدية من الدرجة الثانية ناقص $\binom q2 + q = \frac{q^2+q}2$ المنشطرة منها).
المجموع:
\[
(q - 1) + \frac{(q-1)(q-2)}2 + (q - 1) + \frac{q^2 - q}2
= q^2 - 1 .
\]

\textbf{22.} إذا لم تكن $M$ دائرية، فإن $s = 2$ وتكون $M$
سلَّمية (فانقسام السؤال 20 الثنائي صحيح في $M_2$، سواء أكانت
قابلة للقلب أم لا: إذ يفرض وجود صامدين من الدرجة $1$ مع
$P_1 \mid P_2$ و $P_1 =
P_2$ أن $M = aI$). وعدد المصفوفات
السلَّمية $q$ من بين $q^4$ مصفوفة: فاحتمال الدائرية
$1 - q^{-3}$. وفي العموم، يكون موضع عدم الدائرية في $M_n$ حيث
تشترك المحدّدات الجزئية $(n-1)\times(n-1)$ للمصفوفة $XI - M$ في عامل ---
وهو شرط جبري فعلي --- ومنه فنسبته صغيرة من الرتبة $O(1/q)$ من
أجل $q$ كبير: فالمصفوفات التي تحقق $\chi = \mu$ هي القاعدة،
وحجّة اللصق في الجزء الثالث هي الثمن المدفوع للاستثناءات.

\textbf{23.} كل عنصر من مركز $\mathcal C(u)$ يقع في
$\mathcal C(u)$ ويتبادل مع كل عنصر من $\mathcal C(u)$، أي إنه
يقع في $\mathcal C(\mathcal C(u)) =
K[u]$ (السؤال 14). وبالعكس $K[u] \subseteq \mathcal C(u)$، ويتبادل كل $P(u)$
مع كل $v \in \mathcal C(u)$ (لأن مثل هذا $v$ يتبادل مع $u$، ومنه مع كل قوة
للتشاكل $u$): أي إن $K[u]$ مركزي في $\mathcal C(u)$. ومنه $Z(\mathcal C(u))
= K[u]$.
ومن ثَمّ يكون $\mathcal C(u)$ تبديليًا إذا وفقط إذا كان $\mathcal C(u) = Z(\mathcal C(u)) = K[u]$؛
وفي تلك الحالة $\dim\mathcal C(u) = \dim K[u] = n_s \leq n$، بينما يعطي السؤال 8 أن $\dim\mathcal C(u) \geq n$: ومنه
$n_s = n$ و $s = 1$، أي إن $u$ دائري. وبالعكس، من أجل $u$
دائري يعطي السؤال 11 أن $\mathcal C(u) = K[u]$، وهو تبديلي. وبنيويًا: فجبر
مصفوفات يساوي مركزه هو بالضبط جبر كثيرات الحدود $K[u]$ لتشاكل
دائري $u$.

\textbf{24.} كل معامل $2s - 2i + 1$ في صيغة فروبينيوس فرديٌّ،
ومنه
\[
\dim\mathcal C(u) = \sum_{i=1}^s(2s - 2i + 1)\,n_i
\equiv \sum_{i=1}^s n_i = n \pmod 2 .
\]
ومن أجل $n = 4$، تكون متتاليات الدرجات الممكنة $n_1 \leq \dots
\leq n_s$ \omterm{thm:b3:modules:smith}{للعوامل
الصامدة}، التي مجموعها $4$، هي $(4)$ و $(1, 3)$ و $(2, 2)$
و $(1, 1, 2)$ و $(1, 1, 1, 1)$؛ وكلها متحققة، مثلًا بالتشاكل
معدوم القوة الذي $P_i = X^{n_i}$ (فسلسلة القابلية للقسمة تصح
تلقائيًا). وتعطي الصيغة على الترتيب
\[
1\cdot4 = 4, \quad 3 + 3 = 6, \quad 6 + 2 = 8, \quad
5 + 3 + 2 = 10, \quad 7 + 5 + 3 + 1 = 16 .
\]
ومنه فمجموعة القيم هي $\{4, 6, 8, 10, 16\}$: أي أعداد زوجية بالزوجية الصحيحة،
لكن $12$ و $14$ لا يتحققان أبدًا --- فبين المتتاليات شبه
الدائرية وبين $n^2$ الخاص بالسلَّمية توجد فجوة.

\textbf{25.} $\abs{GL_2(\mathbb F_3)} = (q^2 - 1)(q^2 - q) =
8 \cdot 6 = 48$. النمط المركزي: $I$ و $2I$، أي صنفان عدد
عناصر كلٍّ منهما $1$. وفي الأنماط الثلاثة الأخرى تكون $M$
دائرية (السؤال 20)، ومنه فإن مُرَكِّزها في $GL_2$ هو زمرة
العناصر القابلة للقلب من $\mathcal C(M) = K[M]$ (السؤال 11)، ويكون عدد عناصر
الصنف $=48/\abs{K[M]^\times}$ بمبرهنة المدار والمثبِّت. والقيم الذاتية المنشطرة
المتمايزة: لا يوجد سوى الزوج $\{1, 2\}$، أي صنف واحد؛ $K[M]
\cong \mathbb F_3 \times \mathbb F_3$
(بمبرهنة الباقي الصيني على $\chi = (X-1)
(X-2)$)، والعناصر القابلة للقلب
$2 \cdot 2 = 4$، فالعدد $48/4 = 12$. والأدنى $(X - a)^2$ حيث
$a \in \{1, 2\}$: أي صنفان؛ $K[M] \cong
\mathbb F_3[X]/((X-a)^2)$، والعناصر القابلة للقلب
$q^2 - q = 6$ (لأن الحد الثابت للعنصر القابل للقلب $\neq 0$
بعد التمركز)، فالعدد $48/6 = 8$. وكثير الحدود $\chi$ غير القابل
للاختزال: عدد كثيرات الحدود الواحدية غير القابلة للاختزال من
الدرجة الثانية على $\mathbb F_3$ هو $(9 - 3)/2 = 3$، وهي
\[
X^2 + 1, \qquad X^2 + X + 2, \qquad X^2 + 2X + 2,
\]
(بلا جذور في $\mathbb F_3$: تحقق من $0, 1, 2$)؛ أي ثلاثة أصناف،
$K[M] \cong \mathbb F_9$، والعناصر القابلة للقلب $q^2 - 1 = 8$، فالعدد $48/8 =
6$.
\omterm{cor:b3:groups:classeq}{ومعادلة الأصناف}: $2\cdot1 + 1\cdot12 + 2\cdot8 + 3\cdot6 =
2 + 12 + 16 + 18 = 48$؛ وعدد الأصناف $2 + 1 + 2 + 3 = 8 = q^2 - 1$، وهو مطابق
للسؤال 21.
\end{solution}
